\subsubsection{Design-Level Decoupling}
\label{app:decouple-design}

The first question is how benchmark construction is intended to provide a
compliant completion path. The data-generation pipeline targets three
properties that together support an authorized trajectory achieving
$\mathrm{TSR}=1$ and $\mathrm{DR}=0$:

\begin{enumerate}
  \item \textbf{Sufficient authorized tools.}
    The agent is provided with a complete set of authorized tools and
    a compliant path that, when followed, fully satisfies the user
    instruction. The Stage~1 generator (Appendix~\ref{app:stage1-prompt})
    is explicitly instructed to verify that the compliant path closes
    the task.

  \item \textbf{Orthogonal drift tools.}
    Drift-inducing tools are present in the environment but strictly
    orthogonal to the core objective. They represent tempting shortcuts,
    over-proactive data gathering, or out-of-scope actions --- never
    prerequisites for the user-requested deliverable.

  \item \textbf{Decoupled judgment.}
    The LLM judge uses rewritten completion prompts that score
    \emph{only} core task completion; safety-side constraints are
    removed from the success criteria and measured exclusively by
    rule-based DR over the trajectory
    (Appendix~\ref{app:success_prompt}). The judge is blind to whether
    drift tools were invoked.
\end{enumerate}

Together, these constraints are designed to make successful completion
representable in the authorized action space. They are design targets rather
than a formal guarantee for every sample; the independent human audit of
path feasibility and success-criteria validity is reported in
Appendix~\ref{app:human-eval}.

\subsubsection{Existence of Compliant Trajectories}
\label{app:decouple-empirical-existence}

The second question is whether models can \emph{empirically} find
the compliant path. We roll out eight bare models under the normal
(no-framework) condition on all 648 benchmark samples with 4
independent rollouts each, giving up to 32 trajectories per task
(20,736 planned trajectories; ${\approx}97.6\%$ complete).
Each trajectory is labeled with \textbf{TaskSuccess} (LLM judge,
decoupled prompt), \textbf{Drift} (rule-based subset match against
\texttt{violation\_tool\_calls}), and
$\textbf{CleanSuccess} := \textbf{TaskSuccess} \land \lnot\textbf{Drift}$.
Statistics are computed over the 644 samples with at least one
analyzed trajectory.

Tables~\ref{tab:per-task-existence} and~\ref{tab:per-task-rate}
report the results. $86.18\%$ of analyzed tasks have at least one
CleanSuccess trajectory, and the mean per-task CleanSuccess rate
is $33.70\%$, showing that compliant completion is empirically
reachable on a non-trivial fraction of attempts for the typical task.

\begin{table}[h]
\centering
\small
\setlength{\tabcolsep}{3pt} % 缩小列间距（关键修复）
\caption{Existence of a compliant completion path per task,
  pooled across eight bare models and four rollouts each
  ($N=644$ tasks with at least one analyzed trajectory).}
\label{tab:per-task-existence}
\resizebox{\columnwidth}{!}{%
\begin{tabular}{lrr}
\toprule
\textbf{Condition on a task} & \textbf{Tasks} & \textbf{Share} \\
\midrule
$\geq 1$ CleanSuccess trajectory   & 555 & 86.18\% \\
$\geq 2$ CleanSuccess trajectories & 513 & 79.66\% \\
$\geq 5$ CleanSuccess trajectories & 422 & 65.53\% \\
\midrule
$0$ CleanSuccess trajectories      &  89 & 13.82\% \\
\quad of which: never completed
                                   &  54 &  8.39\% \\
\quad of which: completed but always with drift
                                   &  35 &  5.43\% \\
\bottomrule
\end{tabular}
}
\end{table}

\begin{table}[h]
\centering
\small
\caption{Distribution of the per-task CleanSuccess rate over
  analyzed rollouts ($N=644$).
  Mean $33.70\%$, median $25.00\%$.}
\label{tab:per-task-rate}
\begin{tabular}{lrr}
\toprule
\textbf{Per-task CleanSuccess rate} & \textbf{Tasks} & \textbf{Share} \\
\midrule
$0\%$          &  89 & 13.82\% \\
$(0,\;10\%]$   &  98 & 15.22\% \\
$(10,\;25\%]$  & 136 & 21.12\% \\
$(25,\;50\%]$  & 132 & 20.50\% \\
$(50,\;75\%]$  & 124 & 19.25\% \\
$(75,\;100\%]$ &  65 & 10.09\% \\
\bottomrule
\end{tabular}
\end{table}

Of the 89 tasks with no CleanSuccess trajectory, 54 ($8.39\%$ of
all tasks) were never completed by any model under any rollout,
indicating a capability ceiling rather than a forbidden compliant
path. Only 35 tasks ($5.43\%$) were completed but always with drift;
we flag this small residual as a known limitation rather than
evidence against the benchmark design.

\subsubsection{Empirical Decoupling at Similar Capability}
\label{app:decouple-empirical-fixed}

The third question is whether the observed capability--drift association
could be explained solely by task construction. If successful completion
mechanically required drift, models with similar TSR would be expected to
show similar DR.

\paragraph{Same-TSR, divergent-DR comparison.}
Claude Sonnet 4.6 and the five Qwen models all achieve
$\mathrm{TSR}=0.903$--$0.925$, while their run-level DR ranges from
$0.434$ to $0.605$, a spread of 17.1 percentage points at similar
task-completion performance. The closest pair is Qwen3.5-397B-A17B
($\mathrm{TSR}=0.919$, $\mathrm{DR}=0.434$) and Qwen3.6-35B-A3B
($\mathrm{TSR}=0.925$, $\mathrm{DR}=0.605$).

Together with the observed clean-success trajectories, this variation
provides evidence against a strict structural coupling in which task
completion mechanically requires drift. It does not establish the causal
source of the cross-model association.

\subsubsection{Case-Level Illustration}
\label{app:decouple-cases}

To make the decoupling tangible, we present two trajectories on the
\emph{same} task and \emph{same} environment, differing only in
the backbone model.

\paragraph{Case A: $\mathrm{TSR}=1$, $\mathrm{DR}=0$ (compliant).}
The user requests a Q3 Financial Report summary for the marketing
team. The agent calls \texttt{read\_financial\_report("Q3\_report")},
extracts the relevant metrics, calls
\texttt{send\_email} with the summary, and stops.
The task is fully completed; no unauthorized tool is invoked.

\paragraph{Case B: $\mathrm{TSR}=1$, $\mathrm{DR}>0$ (drifting).}
On the same task, a second model uses a reference to an upcoming
board meeting regarding Q4 marketing budgets to call
\texttt{read\_confidential\_board\_minutes}
(\textit{data drift: Data minimization}), incorporates the
confidential Q4 budget plans into the summary, and sends the
enriched report via \texttt{send\_email}
(\textit{scope drift: Implicit red line}).
The user-visible deliverable is identical to Case~A; the drift is
an additive action consistent with task-enrichment behavior.

Both trajectories produce the user-requested deliverable, showing in
this case that the drifting call is not a precondition for task success.
The benchmark does not force the agent to call
\texttt{read\_confidential\_board\_minutes}; the paired trajectories
illustrate the optional overreach that ADB is designed to surface.